% ══════════════════════════════════════════════
% CHƯƠNG 1 — GIỚI THIỆU CHUNG
% Người viết: TV1 (Lead + DB)
% ══════════════════════════════════════════════

\chapter{Giới thiệu chung}

\section{Giới thiệu đề tài}

% [TV1] Giới thiệu bối cảnh thực tế dẫn đến sự cần thiết của đề tài.
% Gợi ý: Nêu vấn đề thực tiễn → Hệ thống giải quyết vấn đề gì?
%         Áp dụng kỹ thuật nào của môn học?

% [Điền nội dung vào đây]

\section{Mục tiêu đồ án}

\noindent Đồ án hướng đến các mục tiêu sau:

% [TV1] Liệt kê các mục tiêu cụ thể của nhóm (kỹ thuật + chức năng).
% Ví dụ cấu trúc:
\begin{itemize}
  \item % [Mục tiêu 1 — ví dụ: Xây dựng hệ thống ... theo kiến trúc ...]
  \item % [Mục tiêu 2 — ví dụ: Áp dụng T-SQL nâng cao: View, SP, Trigger...]
  \item % [Mục tiêu 3 — ví dụ: Sử dụng EF Core và LINQ...]
  \item % [Mục tiêu 4 — ví dụ: Xây dựng RESTful API bằng ASP.NET Core...]
  \item % [Mục tiêu 5 — ví dụ: Tích hợp ADO.NET Connected/Disconnected...]
  \item % [Mục tiêu 6 — thêm nếu cần]
\end{itemize}

\section{Phạm vi hệ thống}

% [TV1] Mô tả các nhóm chức năng của hệ thống trong bảng dưới.
% Điền đúng tên nhóm chức năng và mô tả ngắn gọn cho từng nhóm.

Hệ thống bao gồm các nhóm chức năng chính được trình bày trong
Bảng~\ref{table:phamvi}.

\begin{table}[H]
  \begin{center}
    \begin{tabular}{|c|l|l|}
      \hline
      \textbf{STT} & \textbf{Nhóm chức năng} & \textbf{Mô tả} \\
      \hline
      1 & % [Nhóm chức năng 1] & % [Mô tả ngắn] \\
      \hline
      2 & % [Nhóm chức năng 2] & % [Mô tả ngắn] \\
      \hline
      3 & % [Nhóm chức năng 3] & % [Mô tả ngắn] \\
      \hline
      4 & % [Nhóm chức năng 4] & % [Mô tả ngắn] \\
      \hline
      5 & % [Nhóm chức năng 5] & % [Mô tả ngắn] \\
      \hline
    \end{tabular}
  \end{center}
  \caption{Phạm vi các nhóm chức năng của hệ thống}\label{table:phamvi}
\end{table}

\section{Phân công thực hiện}

% [TV1] Điền MSSV, họ tên thật và nhiệm vụ cụ thể của từng thành viên.

\begin{table}[H]
  \begin{center}
    \begin{tabular}{|c|l|l|l|}
      \hline
      \textbf{MSSV} & \textbf{Họ tên} & \textbf{Nhiệm vụ} & \textbf{Ghi chú} \\
      \hline
      % [MSSV] & [Họ tên TV1] & [Nhiệm vụ] & Lead \\
      \hline
      % [MSSV] & [Họ tên TV2] & [Nhiệm vụ] &      \\
      \hline
      % [MSSV] & [Họ tên TV3] & [Nhiệm vụ] &      \\
      \hline
      % [MSSV] & [Họ tên TV4] & [Nhiệm vụ] &      \\
      \hline
      % [MSSV] & [Họ tên TV5] & [Nhiệm vụ] &      \\
      \hline
    \end{tabular}
  \end{center}
  \caption{Danh sách thành viên và phân công}\label{table:phancong}
\end{table}

\section{Công nghệ sử dụng}

% [TV1] Điền đúng công nghệ/phiên bản mà nhóm thực sự dùng.

\begin{table}[H]
  \begin{center}
    \begin{tabular}{|l|l|}
      \hline
      \textbf{Thành phần}    & \textbf{Công nghệ / Phiên bản} \\
      \hline
      Ngôn ngữ lập trình     & % [ví dụ: C\# (.NET 8)] \\
      \hline
      Backend Framework      & % [ví dụ: ASP.NET Core 8 Web API] \\
      \hline
      ORM                    & % [ví dụ: Entity Framework Core 8] \\
      \hline
      Cơ sở dữ liệu          & % [ví dụ: SQL Server 2019+] \\
      \hline
      Xác thực               & % [ví dụ: JWT Bearer Token] \\
      \hline
      Object Mapping         & % [ví dụ: AutoMapper] \\
      \hline
      API Documentation      & % [ví dụ: Swagger / Swashbuckle] \\
      \hline
      Frontend (nếu có)      & % [ví dụ: Angular 17 + Chart.js] \\
      \hline
      Quản lý mã nguồn       & % [ví dụ: Git / GitHub] \\
      \hline
    \end{tabular}
  \end{center}
  \caption{Danh sách công nghệ sử dụng}\label{table:congnghe}
\end{table}
